\chapter{Introducción}
%Escribir un texto de un o dos párrafo(s) máximo de 10 líneas con una introducción al capítulo

Sed interdum tortor at purus lobortis dignissim. Suspendisse sem elit, imperdiet a pharetra at, ultricies eget risus. Nullam eleifend posuere tellus ut fringilla. Maecenas quis leo libero. Pellentesque aliquam sapien eu nunc auctor, a consequat lectus molestie. Proin ullamcorper quam magna, ornare tempus justo rutrum ac. Orci varius natoque penatibus et magnis dis parturient montes, nascetur ridiculus mus. Nullam bibendum consequat massa, sed ornare enim faucibus eget. Sed posuere ac sapien sodales congue.  

\section{Motivación}

Lorem ipsum dolor sit amet, consectetur adipiscing elit. Mauris a tincidunt ligula. Quisque quis sem ut ligula venenatis ultricies auctor eu nisi. In lacinia, diam fermentum cursus scelerisque, velit orci tristique velit, in porttitor nisi purus sed libero. Interdum et malesuada fames ac ante ipsum primis in faucibus. Aliquam sagittis mi vel risus maximus, id volutpat justo tempor. Sed malesuada ante in odio euismod semper sed ut libero. Nulla congue a magna quis vestibulum. Praesent cursus arcu a lectus hendrerit maximus. Nulla lobortis, erat facilisis viverra euismod, lectus nibh rutrum turpis, a dapibus orci nisl ut odio. Duis ultricies auctor orci, tristique lobortis felis sollicitudin sed. 

%Que es lo que te ha motivado para realizar la tesis en esta temática y los aspectos más esenciales que pensabas obtener de ella...

%El habla es uno de los primeros medios que emplea el ser humano para comunicarse. Por ello, desde hace algunos años se investiga el procesamiento del habla con el fin de poder realizar diversas tareas a partir de ella como la transcripción y traducción automática, el reconocimiento de voz, las cuales tienen una gran utilidad. En particular, el reconocimiento del habla tiene diversas aplicaciones como el dictado automático, control por comandos y apoyo para discapacitados.

%Dado que en un inicio no se contaba con la capacidad computacional que se tiene al día de hoy, se emplearon métodos complejos para construir un sistema que realice esta tarea. Estos sistemas contaban con tres modelos: acústico, de lenguaje y de pronunciación, los cuales eran difícil de optimizar. 

%Con el incremento de la capacidad de cómputo y por ende, la introducción del Deep Learning, se logran desarrollar sistemas End-to-End que asignan directamente el audio de entrada a la secuencia de palabras que conforman su transcripción. 


%Este punto podrá ser de 4 páginas máximo. Es una de las partes más importantes que introduce al lector en el trabajo en tu tesis, por lo que debe estar muy bien redactado y estructurado.

\section{Objetivo General}

Vivamus id felis sed sapien egestas finibus. Quisque vulputate sit amet ex rhoncus tempus. Curabitur suscipit eget est eu placerat. Sed dignissim purus vitae eleifend pellentesque. In in posuere erat, eu euismod massa. Ut in nisi elit. Morbi in arcu ut enim semper venenatis lobortis ac odio. Donec varius dui sit amet lacinia placerat. Proin eu justo vitae elit fringilla fringilla eu ut metus. 

Específicamente, los objetivos de este trabajo son:

\begin{itemize}
\item[•] Quisque sed tortor ullamcorper, fermentum mauris vitae, porttitor lacus.
\item[•] Proin aliquam elit ac lectus auctor, nec pharetra massa vulputate.
\item[•] Morbi luctus purus vel dui ornare, vel sollicitudin justo bibendum.

\end{itemize}

Orci varius natoque penatibus et magnis dis parturient montes, nascetur ridiculus mus. Donec mattis porttitor nunc, eget blandit neque sagittis vitae:
\begin{itemize}

\item[•] Sed tempus eros eget augue malesuada, faucibus venenatis quam mattis.
\item[•] Suspendisse eget elit aliquam, semper dolor ac, fermentum purus.
\end{itemize}

\section{Estructura del Seminario}
%Para brindar al lector una idea global del contenido de este trabajo, a continuación se detalla una breve descripción de cada capítulo presente en este seminario de tesis.
\begin{itemize}
\item \textbf{Introducción:} \\
Nullam semper est non tellus vestibulum elementum. In rutrum accumsan orci, non porttitor diam ornare id. Morbi venenatis nibh eleifend velit finibus ornare. Sed sapien libero, finibus ac sem sed, pretium elementum velit. Etiam ullamcorper egestas sapien, in auctor augue accumsan eu. Aliquam sollicitudin semper sem, non rutrum metus. Aenean feugiat ante eu sagittis blandit. Mauris fermentum tempus urna a tincidunt. Aliquam lacinia, risus non consectetur rhoncus, nibh massa rhoncus felis, sit amet tincidunt felis felis ac nulla. Ut mattis semper urna ut hendrerit. Vivamus in ex mattis, finibus metus non, vulputate erat. Integer orci ipsum, aliquet id lorem vitae, rhoncus malesuada ligula. Pellentesque tristique odio in mattis accumsan. 
%En este capítulo introductorio se expone la motivación para desarrollar este trabajo de investigación. Asimismo, se presentan los objetivos y una visión general de la organización del documento.

%En este capítulo introductorio se comenta sobre las motivaciones ...

\item \textbf{Estado del Arte:} \\
nterdum et malesuada fames ac ante ipsum primis in faucibus. Nunc sollicitudin, lorem rhoncus sollicitudin fringilla, nibh turpis consectetur nisl, sed iaculis justo massa at risus. Pellentesque volutpat euismod mi, nec gravida justo dictum id. Etiam iaculis tortor dolor, eget egestas eros luctus non. Quisque consectetur egestas arcu vitae blandit. Etiam pretium tortor ullamcorper nibh porta, sit amet egestas justo auctor. Aliquam gravida venenatis metus id iaculis. 
%La base teórica sobre la cual se sustenta el presente trabajo se detalla en este apartado. Se exponen, además, conceptos con respecto a la navegación de un robot móvil, la utilización de sensores y cámaras para identificar el ambiente donde se movilizará el robot así como los métodos para su localización y algoritmos de enrutamiento.
%Este apartado presenta la forma en como se genera el habla así como también se muestran conceptos básicos de Deep Learning. Finalmente, se expone el problema que busca resolver este trabajo.

\item \textbf{Herramientas:} \\
Pellentesque consequat auctor dolor quis viverra. Aenean a felis pulvinar, imperdiet urna ac, fermentum lorem. Integer porta turpis nec ex fringilla, a posuere tortor tristique. Nullam pulvinar augue nec facilisis ultricies. Sed vel fermentum neque. Fusce erat metus, varius et gravida nec, aliquam a ligula. Pellentesque habitant morbi tristique senectus et netus et malesuada fames ac turpis egestas. 
%El capítulo se centra en la descripción de librerías empleadas para el procesamiento de los datos y la creación de las arquitecturas de red.

%En este capítulo se describe el sistema empleado, sus componentes electrónicos y el framework utilizado para la comunicación entre la tarjeta controladora y la microcomputadora. 

\item \textbf{Análisis de Datos de Entorno} \\
%En este capítulo se detalla la base de datos empleada, las transformaciones que se aplicaron a los datos y el aumento de los datos.
Aliquam ex felis, malesuada vel blandit ac, finibus vitae odio. Nulla facilisi. Etiam euismod vehicula risus, id aliquet nibh. Quisque auctor venenatis molestie.  

\item \textbf{Desarrollo de la investigación:} \\
%En este capítulo se detalla el problema que busca resolver el presente seminario y la propuesta de solución.
%En este capítulo se describe la ejecución del 
Nulla vitae elit nec ligula lobortis aliquam. Aliquam lobortis iaculis nibh semper bibendum. Aliquam feugiat mi et velit laoreet, ac blandit turpis blandit. Phasellus posuere, nisl a vestibulum finibus, nisi urna luctus magna, quis elementum nunc risus sit amet sem. Fusce non ligula interdum, mollis quam at, iaculis lorem. Maecenas consequat volutpat nulla ac scelerisque. 

\item \textbf{Análisis de Resultados:} \\
Interdum et malesuada fames ac ante ipsum primis in faucibus. Ut cursus tempor consectetur. Praesent eleifend arcu ac mattis mollis. Morbi at lobortis nunc.
%En este capítulo se muestran los resultados experimentales obtenidos y se comparan con los resultados esperados.
%En este apartado se analiza la exactitud del sistema ASR desarrollado y se realiza una comparación entre los diversos parámetros utilizados.

\item \textbf{Conclusiones y Trabajo Futuro:}\\
Curabitur a blandit lectus. Etiam dui mi, interdum sed gravida ac, faucibus vel justo. Sed tellus erat, rhoncus at nibh sed, convallis finibus leo. Sed commodo nec arcu nec sodales. 
%El capítulo se centra en exponer las conclusiones obtenidas de este trabajo. Adicionalmente, se proponen trabajos a futuro que puedan ser implementados con el fin de mejorar la navegación.
%El capítulo se centra en exponer las conclusiones de acuerdo a los objetivos propuestos para este Seminario de Tesis II. Adicionalmente, se proponen trabajos a futuro que puedan ser implementados con el fin de mejorar el sistema ASR.
%para la implementación de el servidor y la aplicación web del sistema.
\end{itemize}

%NOTA: RECUERDE QUE ES COMO UN LIBRO TODO CAPÍTULO NUEVO COMENZARÁ CON PÁGINA IMPAR NUNCA PAR

\afterpage{\blankpage}


